\section{Bessel rows with inputs varying in depth}
\label{sec:rows}

We shall apply fixed-input operator estimates to rows that read different
members of one input sequence.  The input is allowed bounded total variation
on each of finitely many consecutive depth blocks.  The additional structure
of the rows is spatial: their adjoints have constant coefficients in a
persistent catalog of columns with private coordinates.  We first prove the
result for such rows, and then treat smooth mean-zero rows with fixed
carriers.

Throughout this section $R$ is a bounded dyadic interval of positive length.
Unqualified $L^2$ norms use Lebesgue measure, without normalization.
The inner product of a complex Hilbert space is linear in its first
variable.  A dyadic interval is taken half open.  A finite dyadic partition
of an interval $R$ consists of pairwise disjoint dyadic subintervals whose
union is $R$; a partition refines another if each of its cells is contained
in one cell of the other.  Thus cells of a partition may have different
depths.  Choose an integer $s_R$ and define the spatial depth of a dyadic
cell $I\subset R$ by
\[
 s(I)=s_R+\log_2\frac{|R|}{|I|}.
\]
For an integer $h\ge s_R$, the full depth-$h$ partition consists of all
cells in $R$ of depth $h$.  On a dyadic subroot $R'\subset R$, retain
$s_{R'}=s(R')$, so the spatial depths of its cells are unchanged.

\subsection{The input size and variation bounds}

Let
\[
 \mathsf S=\{s_-,s_-+1,\ldots,s_+\}
\]
be a nonempty finite interval of integers.  Partition $\mathsf S$ into
$W\ge1$ nonempty consecutive blocks $\mathsf S_1,\ldots,\mathsf S_W$,
in that order.  This partition is fixed independently of the spatial
point.  For strongly measurable functions $F_s:R\to\mathcal H$ into a
Hilbert space, the block normalization will mean
\begin{equation}
 \max_{s\in\mathsf S_b}\|F_s(y)\|_{\mathcal H}
 +\sum_{\substack{s,s+1\in\mathsf S_b}}
       \|F_{s+1}(y)-F_s(y)\|_{\mathcal H}
 \le W^{-1/2}
 \quad (1\le b\le W)
 \label{eq:row-input-normalization}
\end{equation}
for almost every $y\in R$.  Since the data are finite, one may take the
same exceptional null set for all blocks.

\begin{lemma}[Size and variation of the depth inputs]
\label{lem:depth-input-size}
Assume Equation~\eqref{eq:row-input-normalization}.  For every measurable
$E\subset R$ of positive measure, every $s\in\mathsf S$, and
$1\le q\le\infty$, one has
\[
 \|F_s\|_{q,E}\le W^{-1/2}\le1,
 \qquad
 \|\mathbf1_EF_s\|_2\le W^{-1/2}|E|^{1/2}\le |E|^{1/2}.
\]
Here the first norm uses normalized measure on $E$, with the usual
essential-supremum interpretation when $q=\infty$.

Let $\mathcal J$ be a finite collection of pairs $[r,t]$ with
$r,t\in\mathsf S$ and $r\le t$.  Its edge overlap is
\[
 \mathfrak m=
 \max_{s_-\le j<s_+}
       \#\{[r,t]\in\mathcal J:r\le j<t\},
\]
where a maximum over an empty set is zero.  Then
\begin{equation}
 \sum_{[r,t]\in\mathcal J}
       \|\mathbf1_E(F_t-F_r)\|_2^2
 \le \mathfrak m\left(1+\frac{4(W-1)}W\right)|E|
 \le 5\mathfrak m|E|.
 \label{eq:row-variation}
\end{equation}
In particular the constant is independent of the number of blocks and
depths.  The same statements hold on a nondecreasing subsequence of the
schedule, with repeated indices allowed and overlap counted in that
subsequence.
\end{lemma}

\begin{proof}
The size statements follow pointwise from
Equation~\eqref{eq:row-input-normalization}.  To prove the variation
estimate, fix a point outside the exceptional null set and put
\[
 d_j=\|F_{j+1}(y)-F_j(y)\|_{\mathcal H},
 \qquad
 v_b=\sum_{j,j+1\in\mathsf S_b}d_j.
\]
Thus $v_b\le W^{-1/2}$.  For an interval $[r,t]$ with both endpoints in
$\mathsf S_b$, let $\delta_{r,t}=\|F_t(y)-F_r(y)\|_{\mathcal H}$.
The triangle inequality gives
$\delta_{r,t}\le\sum_{j=r}^{t-1}d_j\le v_b$.  Hence
\[
 \sum_{\substack{[r,t]\in\mathcal J\\r,t\in\mathsf S_b}}
       \delta_{r,t}^2
 \le v_b
       \sum_{\substack{[r,t]\in\mathcal J\\r,t\in\mathsf S_b}}
          \sum_{j=r}^{t-1}d_j
 \le \mathfrak m v_b^2.
\]
An interval whose endpoints are in different blocks crosses an edge
between consecutive blocks.  Assign it to its first such edge.  At most
$\mathfrak m$ intervals are assigned to each of the $W-1$ edges.  Each
of their differences has norm at most $2W^{-1/2}$, by the two endpoint
size bounds.  The pointwise sum of all the squared differences is
therefore at most
\[
 \mathfrak m\sum_{b=1}^Wv_b^2+
       4\mathfrak m(W-1)W^{-1}
 \le \mathfrak m\left(1+\frac{4(W-1)}W\right).
\]
Integration on $E$ proves Equation~\eqref{eq:row-variation}.

For a nondecreasing list of indices, every repeated index contributes
zero variation.  Between distinct successive indices, the triangle
inequality charges each original edge at most once.  This proves the
same bound on each nonempty restricted block, retaining the original
$W$ block labels and assigning value zero to an empty block.  Moreover, if
intervals in the list have edge overlap $\mathfrak m$, the original
edge between $j$ and $j+1$ is crossed at most $\mathfrak m$ times:
every interval crossing it crosses the single cut in the list between
indices at most $j$ and indices greater than $j$.  The preceding proof
therefore applies verbatim.  Constant extension before or after the
schedule is also permitted by extending the first or last block.
\end{proof}

A scalar input and its lift $F_s=(f_s,0)$ into a space with private
coordinates have the same pointwise norm and variation.  Thus
Lemma~\ref{lem:depth-input-size} applies to the actual inputs used below on
every subroot.

\subsection{A persistent catalog and the maximal estimate}

We use the private-column construction of \cite[Section~3]{H}, allowing
a Hilbert space in place of its physical scalar coordinate.
Fix a dyadic root $R$ and a finite label set $\mathcal V$.
The cardinality of $\mathcal V$ need not be bounded in terms of the
budget below.  For each
$v\in\mathcal V$ choose a dyadic birth cell $J_v\subset R$, a strongly
measurable function $\phi_v:J_v\to\mathcal H_0$ into a Hilbert space,
and a nonzero scalar $\gamma_v$.  In
\[
 \mathcal H=\mathcal H_0\oplus\ell^2(\mathcal V)
\]
define, once and for all,
\[
 \Phi_v(y)=(\phi_v(y),\gamma_v\mathbf e_v),\qquad y\in J_v.
\]
The direction $\mathbf e_v$ is private to $v$, including when two
physical functions coincide.  On a descendant cell $Q$ the available
labels are
\[
 \mathcal V(Q)=\{v\in\mathcal V:Q\subset J_v\}.
\]
Their columns on $Q$ are the restrictions of these same functions
$\Phi_v$.  An indicator times such a restricted column denotes its
extension by zero off that cell.  This is what persistence means here.
A mask may omit a
label at a particular use, but such omission does not remove the label
from the catalog or change its function on later descendants.  We
assume that $D\ge2$ and
\begin{equation}
 \|\Phi_v\|_{L^\infty(J_v;\mathcal H)}\le D,\qquad
 |\gamma_v|\ge D^{-1},\qquad
 \sum_{v\in\mathcal V}\mathbf1_{J_v}(y)\le D
 \quad\hbox{for almost every }y\in R.
 \label{eq:row-catalog-budget}
\end{equation}
The last bound counts all birth labels on a path, including labels not
selected in a current mask.

\begin{lemma}[Bounds for the private catalog]
\label{lem:row-catalog}
For every dyadic cell $Q\subset R$ and every vector
$a=(a_v)_{v\in\mathcal V(Q)}$, one has
\begin{equation}
 D^{-1}\|a\|_{\ell^2}
 \le \left\|\sum_{v\in\mathcal V(Q)}a_v\Phi_v(y)\right\|_{\mathcal H}
 \le D^{3/2}\|a\|_{\ell^2}
 \quad\hbox{for almost every }y\in Q.
 \label{eq:row-catalog-comparison}
\end{equation}
In particular constant coefficients in such a representation are
unique.  Coefficients may be regarded as vectors in the one space
$\ell^2(\mathcal V)$ by extending them by zero; the vector for a
combination on $Q$ is unchanged when the combination is restricted
to a descendant.

These assertions and the same value of $D$ hold after restriction to
a dyadic subroot $R'\subset R$: discard birth cells disjoint from
$R'$, give an inherited label with $R'\subset J_v$ the new birth
cell $R'$, and retain the other birth cells contained in $R'$.
\end{lemma}

\begin{proof}
Outside one common null set, the private component of the sum has
squared norm $\sum_v|\gamma_v|^2|a_v|^2$, which is at least
$D^{-2}\|a\|_{\ell^2}^2$.  Cauchy--Schwarz and the path count give
\[
 \left\|\sum_va_v\Phi_v(y)\right\|_{\mathcal H}
 \le \left(\sum_{v:y\in J_v}\|\Phi_v(y)\|_{\mathcal H}^2\right)^{1/2}
       \|a\|_{\ell^2}
 \le D^{3/2}\|a\|_{\ell^2}.
\]
The null set can be chosen common because $\mathcal V$ is finite.
The lower bound proves uniqueness.  Availability is monotone on
descendants, and restriction changes neither component of a column,
so the assertion about coefficient vectors follows.  Finally, two
dyadic cells that intersect are nested.  The stated change of births
on $R'$ therefore retains precisely the old columns that can occur
there, with their old values and directions.  It cannot increase
any of the three bounds in Equation~\eqref{eq:row-catalog-budget}.
\end{proof}

We now give the row estimate.  Its hypotheses distinguish the
partition supporting a row from the finer partition on which its
adjoint coefficients are observed.  The order of these two
partitions ensures that every stopping decision assigns an entire
row to one set of indices.

\begin{lemma}[Maximal estimate for local catalog rows]
\label{lem:row-maximal}
Use a root and catalog satisfying
Equation~\eqref{eq:row-catalog-budget}.  Let $N\ge1$, and for
$1\le n\le N$ let $\pi_n$ and $\sigma_n$ be finite dyadic partitions
of $R$ such that
\[
 \sigma_n\text{ refines }\pi_n,\qquad
 \pi_{n+1}\text{ refines }\sigma_n\quad(1\le n<N).
\]
Let $\mathcal R_1,\ldots,\mathcal R_N$ be pairwise disjoint finite
row index sets, with union $\mathcal R$.  For every
$\nu\in\mathcal R_n$, let $\mathcal Y_\nu$ be a Hilbert space,
let $Q_\nu\in\pi_n$, and let
$A_\nu:L^2(R;\mathcal H)\to\mathcal Y_\nu$ be bounded and linear.
Assume the following two properties.
\begin{enumerate}
\item The whole row has support $Q_\nu$:
      $A_\nu=A_\nu\mathbf1_{Q_\nu}$.  For every $K\in\sigma_n$
      and every $c\in\mathcal Y_\nu$, there are spatially constant
      coefficients $a_{\nu,K,v}(c)$ such that, in $L^2(R;\mathcal H)$,
      \[
       \mathbf1_K A_\nu^*c
       =\mathbf1_K\sum_{v\in\mathcal V(K)}
                        a_{\nu,K,v}(c)\Phi_v.
      \]
      There is no asserted uniform bound on these coefficients.
\item For some $B\ge0$ the original rows satisfy the joint operator
      estimate
      \[
       \sum_{\nu\in\mathcal R}\|A_\nu z\|_{\mathcal Y_\nu}^2
       \le B^2\|z\|_{L^2(R;\mathcal H)}^2
       \qquad(z\in L^2(R;\mathcal H)).
      \]
\end{enumerate}
Cells with no rows and empty rounds are allowed, as are multiple rows with
the same support cell.  Define
\[
 L_D=1+\lceil5\log_2D\rceil,\qquad
 \mathfrak C_D=\frac83\sqrt{9+L_D^2}
 \le \frac{40}{3}\log_2(2D).
\]
For $1\le p\le q\le N$ and any tests
$c_\nu\in\mathcal Y_\nu$ in these rounds, put
\[
 \begin{split}
 E_c&=\sum_{n=p}^q\sum_{\nu\in\mathcal R_n}\|c_\nu\|_{\mathcal Y_\nu}^2,\\
 g_n&=\sum_{\nu\in\mathcal R_n}A_\nu^*c_\nu,\qquad
 S_j=\sum_{n=p}^j g_n,\quad S_{p-1}=0,\\
 M(y)&=\max_{p-1\le j\le q}\|S_j(y)\|_{\mathcal H}.
 \end{split}
\]
Then, for every $\lambda>0$,
\begin{equation}
 |\{y\in R:M(y)>\lambda\}|
 \le \min\left\{|R|,
       \frac{\mathfrak C_D^2 B^2 E_c}{4\lambda^2}\right\},
 \qquad
 \|M\|_{L^1(R)}
 \le \mathfrak C_D B\,|R|^{1/2}E_c^{1/2}.
 \label{eq:row-maximal}
\end{equation}
\end{lemma}

\begin{proof}
Restrict to the rounds $p,\ldots,q$ and relabel them starting at one.
Fix the test vector $c$ and $\lambda>0$.  Choose representatives of
the finitely many tested adjoints in their asserted cellwise forms,
and remove their exceptional null sets together with the catalog null
set.  On a cell $K\in\sigma_n$, the resulting function $g_n$ has a
unique coefficient vector $a_n(K)\in\ell^2(\mathcal V)$, supported
on $\mathcal V(K)$, by Lemma~\ref{lem:row-catalog}.  The locality
identity also gives
$A_\nu^*=\mathbf1_{Q_\nu}A_\nu^*$.  Thus on $K$ only rows supported
on the unique cell of $\pi_n$ containing $K$ contribute to $a_n(K)$.

Set
\[
 T=2^{\lceil5\log_2D\rceil},\qquad
 \tau=\frac{\lambda}{4D^{3/2}}.
\]
We will select $T$ disjoint subfamilies of whole rows so that, away from
paths with $T$ hits, every partial sum is a prefix of completed bucket
sums plus a remainder of norm at most $\lambda/4$.  The upper catalog
bound controls the remainder; the lower bound and the original Bessel
inequality control the measure of paths with $T$ hits.  A binary-prefix
estimate then costs only $\log T$.

We assign whole rows to buckets $1,\ldots,T$ as follows.  At the
start of the first round, every cell has current bucket $1$ and
current coefficient vector zero.  Inductively these two states at
the start of round $n$ are constant on each cell $Q\in\pi_n$,
and the vector is supported on $\mathcal V(Q)$.
Assign every row with support $Q$ to its current bucket, unless
that bucket is $T+1$, in which case omit the row.  On each
$K\in\sigma_n$ contained in $Q$ with current bucket at most $T$,
add $a_n(K)$ to the current coefficient vector.  If the norm of the resulting vector is
greater than $\tau$, declare a hit, move to the next bucket, and
set the current vector to zero.  Otherwise retain the resulting
vector and the current bucket.  Once bucket $T$ has a hit, keep
state $(T+1,0)$ and make no further additions.

This induction is well defined on whole cells.  Indeed, the new
states are constant on cells of $\sigma_n$.  All their coefficients
refer to the same global label space: older columns and their
coefficients restrict unchanged to $K$, and their labels remain
available there.  Since $\pi_{n+1}$ refines $\sigma_n$, the states
are constant on each support cell at the start of the next round.
In particular a row is assigned one bucket for its entire support,
and is never split by the set on which a hit has occurred.  All
hit sets are finite unions of cells and are measurable.

Figure~\ref{fig:row-chronology} illustrates the refinement that keeps this
assignment on whole rows.
\begin{figure}[ht]
\centering
\begin{tikzpicture}[x=0.92cm,y=1cm,box/.style={draw,minimum height=6mm},every node/.style={font=\small}]
  \node[anchor=east] at (-0.25,2.7) {support cells $\pi_n$};
  \node[anchor=east] at (-0.25,1.5) {observation cells $\sigma_n$};
  \node[anchor=east] at (-0.25,0.3) {support cells $\pi_{n+1}$};

  \draw[fill=black!3] (0,2.4) rectangle (8,3.0);
  \node at (4,2.7) {one row on $Q$, assigned to bucket $b$};

  \draw[fill=black!3] (0,1.2) rectangle (4,1.8);
  \draw[fill=black!13] (4,1.2) rectangle (8,1.8);
  \node at (2,1.5) {no new hit};
  \node at (6,1.5) {a new hit};

  \draw[fill=black!3] (0,0) rectangle (2,0.6);
  \draw[fill=black!3] (2,0) rectangle (4,0.6);
  \draw[fill=black!13] (4,0) rectangle (6,0.6);
  \draw[fill=black!13] (6,0) rectangle (8,0.6);
  \node at (1,0.3) {$b$};
  \node at (3,0.3) {$b$};
  \node at (5,0.3) {$b+1$};
  \node at (7,0.3) {$b+1$};

  \draw[densely dashed,black!35] (0,0.6)--(0,1.2) (4,0.6)--(4,1.2) (8,0.6)--(8,1.2);
  \draw[densely dashed,black!35] (0,1.8)--(0,2.4) (8,1.8)--(8,2.4);
  \draw[-{Latex[length=2mm]},black!65] (8.45,2.7)--(8.45,0.3);
  \node[anchor=west,align=left] at (8.6,1.5) {chronological\\order};
\end{tikzpicture}

\caption{Schematic view of one branch of the spatial partitions in consecutive
rounds.  The row on $Q$ has already entered bucket $b$ before a new hit is tested.
Hits may differ on the observation cells.  Each support cell in the next
round lies inside one observation cell, so every row in that round still
receives one bucket; the bottom labels give those bucket numbers.}
\label{fig:row-chronology}
\end{figure}

Let $\mathcal B_i$ be the rows assigned to bucket $i$, and set
\[
 v_i=\sum_{\nu\in\mathcal B_i} A_\nu^*c_\nu,\qquad 1\le i\le T.
\]
For every subset $J\subset\{1,\ldots,T\}$, the adjoint of the
joint operator in the hypothesis gives
\begin{equation}
 \left\|\sum_{i\in J}v_i\right\|_2^2
 \le B^2\sum_{\nu\in\bigcup_{i\in J}\mathcal B_i}
                   \|c_\nu\|_{\mathcal Y_\nu}^2.
 \label{eq:row-bucket-bessel}
\end{equation}
Although the sets $\mathcal B_i$ depend on the fixed test vector,
they are actual subsets of the original row indices.  Thus
Equation~\eqref{eq:row-bucket-bessel} is the original operator
inequality applied to a test vector with some coordinates set to
zero.  In particular
\[
 \sum_{i=1}^T\|v_i\|_2^2\le B^2E_c.
\]

Let $\Omega_T$ be the set of points whose path has $T$ hits during
these rounds.  At any point $y$, a tested adjoint from $\mathcal B_i$
whose support does not contain $y$ vanishes there.  If its support
contains $y$, it was assigned $i$ exactly when the path's current
bucket was $i$; conversely, before stopping, every tested adjoint
whose support contains $y$ enters that current bucket.
Thus, on $\Omega_T$, $v_i(y)$ is exactly the sum accumulated in
bucket $i$ up to and including the round that causes its hit.
The coefficient vector of this sum has norm greater than $\tau$.  All its labels remain
available on the final observation cell containing $y$, so
Equation~\eqref{eq:row-catalog-comparison} yields
$\|v_i(y)\|_{\mathcal H}>D^{-1}\tau$ for $1\le i\le T$.
Consequently
\begin{equation}
 |\Omega_T|
 \le \frac{B^2E_c}{TD^{-2}\tau^2}
 =\frac{16D^5B^2E_c}{T\lambda^2}.
 \label{eq:row-many-hits}
\end{equation}

Outside $\Omega_T$, at any intermediate round the partial sum
equals a prefix of the completed bucket sums $v_i(y)$ plus the
current unfinished sum.  The latter has coefficient norm at most
$\tau$, and hence has Hilbert norm at most
$D^{3/2}\tau=\lambda/4$.  This statement also includes the instant
of a hit, by putting the newly completed bucket in the prefix and
taking unfinished sum zero.  Thus, with
\[
 M_T(y)=\max_{0\le j\le T}\left\|\sum_{i=1}^jv_i(y)\right\|_{\mathcal H},
\]
one has
\begin{equation}
 M(y)\le M_T(y)+\lambda/4\qquad(y\notin\Omega_T)
 \label{eq:row-prefix-domination}
\end{equation}
outside the fixed null set.

It remains to estimate $M_T$ using
Equation~\eqref{eq:row-bucket-bessel}.  We use the binary-prefix argument
of Rademacher--Menshov type; compare \cite[Lemma~2.7]{KrauseLacey2020}.
For each
$0\le k\le\log_2T$, let $\mathcal D_k$ be the partition of
$\{1,\ldots,T\}$ into consecutive dyadic index intervals of length
$2^k$, and put $v_J=\sum_{i\in J}v_i$.  Every prefix is the disjoint
union of at most one interval from each $\mathcal D_k$.  The triangle
inequality therefore gives pointwise
\[
 M_T\le
 \sum_{k=0}^{\log_2T}
       \left(\sum_{J\in\mathcal D_k}\|v_J\|_{\mathcal H}^2\right)^{1/2}.
\]
For a fixed $k$, Equation~\eqref{eq:row-bucket-bessel} and the
disjointness of the index intervals imply
\[
 \int_R\sum_{J\in\mathcal D_k}\|v_J(y)\|_{\mathcal H}^2\,dy
 =\sum_{J\in\mathcal D_k}\|v_J\|_2^2
 \le B^2E_c.
\]
Minkowski's inequality in $L^2(R)$ now gives
\begin{equation}
 \|M_T\|_2\le (1+\log_2T)B E_c^{1/2}=L_D B E_c^{1/2}.
 \label{eq:row-binary-prefix}
\end{equation}

Equations~\eqref{eq:row-many-hits}--\eqref{eq:row-binary-prefix}
and Chebyshev's inequality show that
\[
 \begin{split}
 |\{M>\lambda\}|
 &\le |\Omega_T|+|\{M_T>3\lambda/4\}|\\
 &\le \frac{16B^2E_c}{\lambda^2}
             \left(\frac{D^5}{T}+\frac{L_D^2}{9}\right)
 \le \frac{16(9+L_D^2)B^2E_c}{9\lambda^2},
 \end{split}
\]
because $T\ge D^5$.  This is the stated weak estimate, as
$\mathfrak C_D^2/4=16(9+L_D^2)/9$.  Integrating the minimum of
this estimate and $|R|$ over $\lambda>0$ proves the $L^1$ estimate:
for $H=(\mathfrak C_D/2)BE_c^{1/2}>0$,
\[
 \int_0^\infty\min\{|R|,H^2/\lambda^2\}\,d\lambda
 =2H|R|^{1/2}.
\]
If $H=0$, the weak estimate gives $M=0$ almost everywhere instead.
Finally, put $x=\log_2D\ge1$.  Since $L_D\le5x+2$ and
$9+(5x+2)^2\le25(x+1)^2$, the asserted upper bound on
$\mathfrak C_D$ follows.
\end{proof}

\begin{lemma}[The depth-varying row estimate]
\label{lem:row-depth}
Let the root, catalog, partitions, and rows satisfy every hypothesis
of Lemma~\ref{lem:row-maximal}.  Let $(F_s)_{s\in\mathsf S}$ satisfy
Equation~\eqref{eq:row-input-normalization} on $R$.  Choose
$s_n\in\mathsf S$ with
$s_1\le s_2\le\cdots\le s_N$, and let every row of round $n$
read the same input $F_{s_n}$.  These read indices label the input
schedule independently of the spatial depths of the support cells.  Then
\begin{equation}
 \left(\sum_{n=1}^N\sum_{\nu\in\mathcal R_n}
          \|A_\nu F_{s_n}\|_{\mathcal Y_\nu}^2\right)^{1/2}
 \le \mathfrak C_D B\,|R|^{1/2}
 \le \frac{40}{3} B\log_2(2D)\,|R|^{1/2}.
 \label{eq:row-depth}
\end{equation}
The constant is independent of the number of rows, depths, and
input blocks.
\end{lemma}

\begin{proof}
Test the output vector against arbitrary
$c=(c_\nu)_{\nu\in\mathcal R}$ in the finite Hilbert direct sum of
the row spaces.  The rounds whose read index lies in a particular
block $\mathsf S_b$ are consecutive, or absent, because the
$s_n$ are nondecreasing.  On a nonempty such range $p\le n\le q$,
form $g_n,S_n,M$ as in Lemma~\ref{lem:row-maximal}.  Pointwise Abel
summation, with the stipulated inner-product convention, is
\[
 \sum_{n=p}^q\langle F_{s_n},g_n\rangle_{\mathcal H}
 =\langle F_{s_q},S_q\rangle_{\mathcal H}
   +\sum_{n=p}^{q-1}
       \langle F_{s_n}-F_{s_{n+1}},S_n\rangle_{\mathcal H}.
\]
The absolute value of this expression is at most
\[
 \left(\|F_{s_q}\|_{\mathcal H}
       +\sum_{n=p}^{q-1}\|F_{s_{n+1}}-F_{s_n}\|_{\mathcal H}\right)M
 \le W^{-1/2}M.
\]
For the last inequality, the nondecreasing reads traverse each
edge in the block at most once; repeated reads contribute zero.
Integrating and using Equation~\eqref{eq:row-maximal} bounds this
block's pairing by
\[
 \mathfrak C_D B\,W^{-1/2}|R|^{1/2}
       \left(\sum_{n=p}^q\sum_{\nu\in\mathcal R_n}
                         \|c_\nu\|_{\mathcal Y_\nu}^2\right)^{1/2}.
\]
All sums are finite and each pairing is integrable by the $L^2$
inequality, so the pointwise identity may be integrated.  Write
$E_b$ for the squared coefficient sum in the last display, with
$E_b=0$ for an absent block.  The adjoint identity and
Cauchy--Schwarz over the $W$ blocks give
\[
 \begin{split}
 \left|\sum_{n=1}^N\sum_{\nu\in\mathcal R_n}
              \langle A_\nu F_{s_n},c_\nu\rangle_{\mathcal Y_\nu}\right|
 &\le \mathfrak C_D B\,|R|^{1/2}
                  W^{-1/2}\sum_{b=1}^W E_b^{1/2}\\
 &\le \mathfrak C_D B\,|R|^{1/2}
        \left(\sum_{\nu\in\mathcal R}\|c_\nu\|_{\mathcal Y_\nu}^2\right)^{1/2}.
 \end{split}
\]
Duality in the finite Hilbert direct sum proves
Equation~\eqref{eq:row-depth}.
\end{proof}

The realized finite rows, partitions, and catalog may depend on the given
depth inputs, provided they satisfy all preceding hypotheses and are held
fixed while their ordinary operator inequality is tested on arbitrary inputs.

\subsection{Projection increments}

Here is a useful consequence for the inherited local projections
of H.  It also records why their splitting into support cells
preserves the original Bessel bound.

\begin{corollary}[Increments of inherited projections]
\label{cor:row-projections}
Use the catalog of Lemma~\ref{lem:row-catalog}.  On a finite
interval of integers $h\ge s_R$, let $\mathcal P_h$ be the full
partition of $R$ into cells of depth $h$.  On each
$Q\in\mathcal P_h$, let $P_{h,Q}$ be the orthogonal projection in
$L^2(Q;\mathcal H)$ onto constant-coefficient combinations of
a subset $\mathcal M_h(Q)\subset\mathcal V(Q)$.  Assume inheritance:
if $Q'\in\mathcal P_{h+1}$ is contained in $Q$, then
$\mathcal M_h(Q)\subset\mathcal M_{h+1}(Q')$.  Let $P_h$ be the
direct sum of the $P_{h,Q}$ on $L^2(R;\mathcal H)$.

Suppose $(F_s)_{s\in\mathsf S}$ satisfies
Equation~\eqref{eq:row-input-normalization}.  For any integer $J\ge1$ and
\[
 r_1\le t_1\le r_2\le t_2\le\cdots\le r_J\le t_J
\]
in the projection depth range and nondecreasing
$s_1,\ldots,s_J\in\mathsf S$, one has
\begin{equation}
 \sum_{j=1}^J\|(P_{t_j}-P_{r_j})F_{s_j}\|_2^2
 \le \mathfrak C_D^2|R|.
 \label{eq:row-disjoint-increments}
\end{equation}
If the intervals can be divided into $\mathfrak m$ such ordered
families whose reads are nondecreasing, the right side is multiplied
by $\mathfrak m$.

For each $s\in\mathsf S$ with $s\ge s_R$, let $\mathcal U_s$ be any
subcollection of the depth-$s$ cells of $R$.  In particular fix an integer
$\ell$ and retain only those $s$ for which $s+\ell$ and $s+\ell+1$ are in
the projection range.  Then
\begin{equation}
 \sum_s\sum_{I\in\mathcal U_s}|I|\,
       \|(P_{s+\ell+1}-P_{s+\ell})F_s\|_{2,I}^2
 \le \mathfrak C_D^2|R|.
 \label{eq:row-increment}
\end{equation}
For every integer $L\ge0$, retaining only $s\in\mathsf S$ with $s\ge s_R$
for which $s$ and $s+L$ are in the projection range,
\begin{equation}
 \left(\sum_s\sum_{I\in\mathcal U_s}|I|\,
       \|(P_{s+L}-P_s)F_s\|_{2,I}^2\right)^{1/2}
 \le L\mathfrak C_D|R|^{1/2}
 \le (1+L)\mathfrak C_D|R|^{1/2}.
 \label{eq:row-delay}
\end{equation}
\end{corollary}

\begin{proof}
The range of $P_h$ is contained in the range of $P_{h+1}$:
restrict a constant-coefficient combination on $Q$ to its
children and use inheritance.  Thus the $P_h$ are nested
orthogonal projections.  For $r\le t$, $P_t-P_r$ is an
orthogonal projection, and the projections corresponding to
intervals with disjoint interiors have pairwise orthogonal ranges.

For a cell $Q\in\mathcal P_r$, both $P_r$ and $P_t$ commute with
multiplication by $\mathbf1_Q$, since both of their defining
partitions refine $\mathcal P_r$.  The row
\[
 A_{j,Q}=(P_{t_j}-P_{r_j})\mathbf1_Q:
        L^2(R;\mathcal H)\longrightarrow L^2(R;\mathcal H),
 \qquad Q\in\mathcal P_{r_j},
\]
has both input and output supported in $Q$.  Consequently
\[
 \sum_{Q\in\mathcal P_{r_j}}\|A_{j,Q}z\|_2^2
 =\|(P_{t_j}-P_{r_j})z\|_2^2,
 \qquad
 \sum_{j,Q}\|A_{j,Q}z\|_2^2\le\|z\|_2^2.
\]
This establishes the joint Bessel bound after the splitting.
The adjoint of $A_{j,Q}$ is
$\mathbf1_Q(P_{t_j}-P_{r_j})$.  On a cell of
$\mathcal P_{t_j}$ its first term is a constant-coefficient
combination there, while its second term is the restriction of
such a combination on an $r_j$-cell.  Both use labels available
on the $t_j$-cell.  Therefore the support and observation
partitions in Lemma~\ref{lem:row-depth} are
$\mathcal P_{r_j}$ and $\mathcal P_{t_j}$, respectively.
The stated order of intervals gives the required refinement,
and Lemma~\ref{lem:row-depth} with $B=1$ proves
Equation~\eqref{eq:row-disjoint-increments}.  Summing the
estimates for $\mathfrak m$ subfamilies proves the overlap version.
After discarding the zero increments $r=t$, intervals with edge overlap
$\mathfrak m$ can be so divided by the usual greedy coloring in increasing left
endpoint, whenever the reads in that order are nondecreasing:
at a left endpoint fewer than $\mathfrak m$ colors are occupied
by earlier intervals whose right endpoints are still larger.

For fixed $\ell$, the intervals $[s+\ell,s+\ell+1]$ are
consecutive and their reads $s$ increase.  Apply
Equation~\eqref{eq:row-disjoint-increments}.  At each $s$ the
cells in $\mathcal U_s$ are disjoint, and
$|I|\|z\|_{2,I}^2=\|\mathbf1_Iz\|_2^2$.  Restricting the
outputs to those cells can only decrease their squared sum,
which proves Equation~\eqref{eq:row-increment}.  Finally,
\[
 (P_{s+L}-P_s)F_s
 =\sum_{\ell=0}^{L-1}(P_{s+\ell+1}-P_{s+\ell})F_s.
\]
Minkowski's inequality in the Hilbert direct sum over $(s,I)$
and Equation~\eqref{eq:row-increment} prove
Equation~\eqref{eq:row-delay}.  For $L=0$ both sides of its
first inequality are zero.
\end{proof}

\subsection{Smooth rows with fixed carriers}

The remaining rows have smooth mean-zero test functions on their
support cells, multiplied by a fixed function that may oscillate
arbitrarily.  Dyadic differences make their coefficients constant
on finer cells.  Private coordinates can then be adjoined at a
controlled cost.

\begin{lemma}[Smooth mean-zero rows on depth inputs]
\label{lem:smooth-depth}
Let $R$ be a finite dyadic root and let $\mathcal N$ be a finite
set of carrier labels.  For each $\nu\in\mathcal N$ fix one
measurable function $m_\nu:R\to\mathbb C$ with
$\|m_\nu\|_\infty\le1$.  This function is the same for every
interval using the label $\nu$.  Let $\mathcal E$ be a finite
set of pairs $(I,\nu)$, where $I\subset R$ is a dyadic interval
of depth $s(I)$ and $\nu\in\mathcal N$.  Thus at most one row
is selected for a given pair $(I,\nu)$.

For each selected pair, let $\omega_{I,\nu}\in C^1(\overline I)$
satisfy
\[
 \int_I\omega_{I,\nu}(y)\,dy=0,\qquad
 \|\omega_{I,\nu}\|_\infty+
       |I|\|\omega_{I,\nu}'\|_\infty\le A
\]
for one $A\ge0$.  Suppose the selected physical rows
\[
 \mathcal W_{I,\nu}f
   =|I|^{-1/2}\int_I m_\nu(y)\omega_{I,\nu}(y)f(y)\,dy
\]
satisfy, for one $B_H\ge0$,
\[
 \sum_{(I,\nu)\in\mathcal E}|\mathcal W_{I,\nu}f|^2
 \le B_H^2\|f\|_{L^2(R)}^2\qquad(f\in L^2(R)).
\]
Let the scalar functions $(f_s)_{s\in\mathsf S}$ satisfy
Equation~\eqref{eq:row-input-normalization}, with absolute
values in place of Hilbert norms, and assume that every
$s(I)$ belongs to $\mathsf S$.  Put
\[
 d=2+\#\mathcal N,\qquad h_0=\lceil10\log_2(2d)\rceil.
\]
Then
\begin{equation}
 \begin{split}
 \left(\sum_{(I,\nu)\in\mathcal E}
           |\mathcal W_{I,\nu}f_{s(I)}|^2\right)^{1/2}
 &\le \mathfrak C_d\left[
       (h_0+1)(B_H+3A)
       +2A\sqrt d\,(h_0+3)2^{-h_0}\right]|R|^{1/2}\\
 &\le 760(B_H+A)\bigl(\log_2(2d)\bigr)^2|R|^{1/2}.
 \end{split}
 \label{eq:smooth-depth}
\end{equation}
The estimate uses only the displayed Bessel inequality for the
selected set $\mathcal E$, not for any larger array of intervals
and carriers.
\end{lemma}

\begin{proof}
The assertion is immediate if $\mathcal E$ is empty.  Extend each
$\omega_{I,\nu}$ by zero off $I$.  Let $\mathbb E_h$ be dyadic
conditional expectation on the full depth-$h$ partition of $R$;
we use it for every integer $h\ge s_R$.  These are nested orthogonal
projections in $L^2(R)$.  For $s=s(I)$, exact mean zero and the zero
extension give $\mathbb E_s\overline{\omega_{I,\nu}}=0$.  Define the
difference pieces for $h\ge1$ and the initial sum by
\[
 \begin{aligned}
 q_{I,\nu,h}
 &=|I|^{-1/2}(\mathbb E_{s+h}-\mathbb E_{s+h-1})
                         \overline{\omega_{I,\nu}}\quad(h\ge1),\\
 q_{I,\nu,0}
 &=|I|^{-1/2}\mathbb E_{s+h_0}\overline{\omega_{I,\nu}}
 =\sum_{h=1}^{h_0}q_{I,\nu,h}.
 \end{aligned}
\]
Every $q_{I,\nu,h}$ is supported in $I$, and for $h\ge1$ it
is constant on depth-$(s+h)$ cells and lies in the martingale
difference space
$\operatorname{range}(\mathbb E_{s+h}-\mathbb E_{s+h-1})$.
The initial sum is constant on depth-$(s+h_0)$ cells.
The zero extension causes no difficulty: cells at these depths
do not cross the boundary of $I$.

On a child of a depth-$(s+h-1)$ cell, the difference between
its average of $\omega_{I,\nu}$ and the parent's average is
at most the oscillation on the parent.  The derivative bound
therefore gives
\begin{equation}
 \|q_{I,\nu,h}\|_2\le 2A\,2^{-h}\qquad(h\ge1).
 \label{eq:smooth-difference-size}
\end{equation}
For a fixed $\nu$ and $h$, these functions are pairwise
orthogonal as $I$ varies through the selected pairs.  Indeed,
different support depths give different martingale difference
spaces.  At one depth the distinct intervals are disjoint,
and there is at most one selected row for $(I,\nu)$.
Thus their analysis operator has norm at most $2A2^{-h}$.
Since multiplication by the one fixed $m_\nu$ is a contraction,
the physical lag-$h$ rows, whose adjoints are
$q_{I,\nu,h}\overline{m_\nu}$, have joint operator norm at most
\begin{equation}
 B_{\mathrm{phys},h}\le 2A\sqrt{\#\mathcal N}\,2^{-h}
 \le 2A\sqrt d\,2^{-h}.
 \label{eq:smooth-physical-lag}
\end{equation}
This argument uses the fixedness and boundedness of $m_\nu$,
and requires no regularity or frequency separation for it.

The derivative bound implies uniform convergence on $I$
of $\mathbb E_{s+H}\omega_{I,\nu}$ to $\omega_{I,\nu}$ as
$H\to\infty$.  Consequently
\[
 |I|^{-1/2}\overline{\omega_{I,\nu}}
 =q_{I,\nu,0}+\sum_{h>h_0}q_{I,\nu,h}.
\]
The associated physical row series also converges in operator norm,
by Equation~\eqref{eq:smooth-physical-lag}.
Subtracting its tail from the original selected row operator
therefore proves
\[
 B_{\mathrm{phys},0}
 \le B_H+2A\sqrt{\#\mathcal N}\,2^{-h_0}
 \le B_H+A.
\]
For the last inequality use
$2^{-h_0}\le(2d)^{-10}$ and $d\ge2$.

Adjoin one private direction for each complete carrier.
Specifically take
$\mathcal H^{\mathrm{aug}}=\mathbb C\oplus\ell^2(\mathcal N)$
and the root-born columns
\[
 \Phi_\nu(y)=(\overline{m_\nu(y)},\mathbf e_\nu),\qquad y\in R.
\]
Their norms are at most $\sqrt2$, their private constants have
size one, and the number of labels on any path is
$\#\mathcal N\le d$.  They satisfy
Equation~\eqref{eq:row-catalog-budget} with $D=d$.
For $h=0$ and for $h>h_0$, define the augmented scalar row by
\[
\widetilde{\mathcal W}_{I,\nu,h}Z
 =\langle Z,q_{I,\nu,h}\Phi_\nu\rangle_{L^2(R;\mathcal H^{\mathrm{aug}})}.
\]
Write $\widetilde{\mathcal W}_h$ for the joint operator consisting
of these rows over $(I,\nu)\in\mathcal E$.
The adjoint of the individual row sends $c$ to
$c\,q_{I,\nu,h}\Phi_\nu$.  On a physical input $Z=(f,0)$,
\[
 \widetilde{\mathcal W}_{I,\nu,h}(f,0)
 =\langle m_\nu f,q_{I,\nu,h}\rangle_{L^2(R)},
\]
which is exactly the physical row for the corresponding bump piece.

For a private input $z=(z_\nu)_{\nu\in\mathcal N}$, the lag-$h$
analysis map has entries $\langle z_\nu,q_{I,\nu,h}\rangle$.
The fixed-$\nu$ orthogonality already proved, followed by
orthogonality of the private coordinates, gives its norm at
most $2A2^{-h}$.  The physical and private input summands are
orthogonal.  Applying Cauchy--Schwarz to their two norm bounds
and using Equation~\eqref{eq:smooth-physical-lag} gives
\begin{equation}
 \|\widetilde{\mathcal W}_h\|
 \le 2A\sqrt{\#\mathcal N+1}\,2^{-h}
 \le 2A\sqrt d\,2^{-h}\qquad(h>h_0).
 \label{eq:smooth-augmented-lag}
\end{equation}
For the initial sum, the identity
$q_{I,\nu,0}=\sum_{h=1}^{h_0}q_{I,\nu,h}$ and the operator
triangle inequality bound its private part by
$2A\sum_{h=1}^{h_0}2^{-h}\le2A$.  Its physical part has the
bound $B_H+A$ above.  Thus
\begin{equation}
 \|\widetilde{\mathcal W}_0\|\le B_H+3A.
 \label{eq:smooth-augmented-main}
\end{equation}

The two pieces, each indexed by the same selected set $\mathcal E$, have
the following data:
\begin{center}
\small
\begin{tabular}{@{}lccc@{}}
\hline
Piece & Observation depth & \shortstack{Joint augmented\\operator norm}
& \shortstack{Residue modulus\\for support depths}\\
\hline
Initial sum $q_0$ & $s+h_0$ & $B_H+3A$ & $h_0+1$\\
Lag $q_h$, $h>h_0$ & $s+h$ & $2A\sqrt d\,2^{-h}$ & $h+1$\\
\hline
\end{tabular}
\end{center}
For either piece, let $\ell$ denote its observation offset, namely $h_0$
or $h$.  Each row is supported on its whole cell $I$ of depth $s=s(I)$,
and its adjoint coefficient is constant on depth-$(s+\ell)$ cells.
Divide the selected support depths into residue classes modulo $\ell+1$.
Within one class, group all selected rows of one depth $s$ in one round,
using the full depth-$s$ support partition and the full depth-$(s+\ell)$
observation partition of $R$; unused cells carry no rows.  Successive
support depths differ by at least $\ell+1$, so the next support partition
refines the preceding observation partition.  All rows in a round read
the same $F_s=(f_s,0)$, and these indices increase in $\mathsf S$.
The catalog has $D=d$, and restricting either displayed array to a
residue class preserves its operator bound.  Lemma~\ref{lem:row-depth}
therefore applies in each of the at most $\ell+1$ classes.

Write $V_0\in\ell^2(\mathcal E)$ for the outputs of the initial sum on
their respective $F_{s(I)}$, and $V_h\in\ell^2(\mathcal E)$ for the
analogous lag-$h$ vector.  Extending each residue-class vector by zero
to the full selected row space and using its triangle inequality yields
\[
 \|V_0\|_{\ell^2(\mathcal E)}
 \le (h_0+1)\mathfrak C_d(B_H+3A)|R|^{1/2},
 \qquad
 \|V_h\|_{\ell^2(\mathcal E)}
 \le (h+1)\mathfrak C_d\,2A\sqrt d\,2^{-h}|R|^{1/2}.
\]
These bounds are summable over $h>h_0$.  The uniform bump
expansion gives rowwise
$\mathcal W_{I,\nu}f_{s(I)}
= (V_0)_{I,\nu}+\sum_{h>h_0}(V_h)_{I,\nu}$.
There are finitely many selected rows; alternatively the last
summable bounds give convergence directly in their Hilbert
sum.  Minkowski's inequality and
\[
 \sum_{h>h_0}(h+1)2^{-h}=(h_0+3)2^{-h_0}
\]
prove the first inequality in Equation~\eqref{eq:smooth-depth}.
For its last inequality put $l=\log_2(2d)\ge2$.  Then
$h_0+1\le11l$, $h_0+3\le12l$, and
$\sqrt d\,2^{-h_0}\le1$.  The bracket there is at most
$57(B_H+A)l$, while
$\mathfrak C_d\le(40/3)l$.  Their product is at most
$760(B_H+A)l^2$, as asserted.
\end{proof}
